\begin{table}[htbp]
    \centering
    \caption{Matriks Keputusan Pemilihan Mekanisme Sinkronisasi Data}
    \label{tab:matriks_sync_engine}
    \begin{tabularx}{\textwidth}{|l|X|X|}
        \hline
        \textbf{Solusi} & \textbf{Kelebihan} & \textbf{Kekurangan} \\
        \hline
        RabbitMQ & Mendukung mekanisme \textit{retry} dan sistem \textit{routing} pesan yang fleksibel, sangat efektif untuk menangani proses asinkron yang membutuhkan keandalan pengiriman data. & Memerlukan instalasi dan pengelolaan \textit{service message broker} tambahan, serta menambah kompleksitas arsitektur dibandingkan solusi bawaan sistem operasi. \\
        \hline
        Apache Kafka & Dirancang untuk \textit{streaming} data bervolume besar dengan tingkat \textit{throughput} yang sangat tinggi serta memiliki toleransi kesalahan (\textit{fault tolerance}) yang superior. & Prosedur instalasi dan pemeliharaan infrastruktur sangat kompleks serta membutuhkan alokasi sumber daya server yang signifikan. \\
        \hline
        Cron Job & Sangat efisien karena merupakan fitur bawaan sistem operasi, ringan, dan tidak memerlukan dependensi eksternal. Sangat ideal untuk kebutuhan \textit{batch processing} atau sinkronisasi berkala. & Tidak bersifat \textit{real-time} (bergantung pada interval penjadwalan) dan memiliki keterbatasan dalam pemantauan otomatis jika terjadi kegagalan proses tanpa integrasi \textit{tooling} tambahan. \\
        \hline
    \end{tabularx}
\end{table}